% Generated draft: requires mathematical and source review.
\documentclass[11pt]{article}
\usepackage[margin=25mm]{geometry}
\usepackage{fontspec}
\setmainfont{Latin Modern Roman}
\usepackage{amsmath,amssymb,mathtools,unicode-math}
\setmathfont{Latin Modern Math}
\usepackage{enumitem,needspace,xurl,hyperref}
\setlength{\parindent}{0pt}
\setlength{\parskip}{5pt}
\setlength{\emergencystretch}{3em}
\newcommand{\N}{\mathbb{N}}
\newcommand{\Z}{\mathbb{Z}}
\newcommand{\Q}{\mathbb{Q}}
\newcommand{\R}{\mathbb{R}}
\newcommand{\C}{\mathbb{C}}
\newcommand{\F}{\mathbb{F}}
\newcommand{\A}{\mathbb{A}}
\newcommand{\PP}{\mathbb{P}}
\newcommand{\E}{\mathbb{E}}
\newcommand{\eps}{\varepsilon}
\newcommand{\dd}{\,\mathrm{d}}
\begin{document}
\section*{UnsolvedMath ID: 30006032}
\textbf{Source problem code:} catalogue entry
\subsection{Definitions and mathematical statement}
The supplied record concerns extending categorical nerve constructions and their universal properties beyond dimension \(2\). A precise formulation would require a category \(\mathcal C\) of higher categorical structures, a category \(\mathcal D\) of target objects, and a nerve functor \(N\colon\mathcal C\to\mathcal D\). It would also require a specified notion of equivalence between nerve constructions and an explicit universal property. These data are not supplied; this notation records the missing ingredients rather than selecting a particular model.

The source excerpt attributes to \([\mathrm{MOR22}]\) a theorem asserting that all \(2\)-categorical nerves are equivalent and satisfy a universal property attributed to Gepner--Haugseng. The excerpt does not specify the class of nerves covered, their domains and targets, the meaning of equivalence, or the universal property. Consequently, the word ``all'' must retain the unspecified scope of that cited result and cannot be read as quantifying over every possible nerve construction.

\emph{Research question, with unresolved scope.} How can the categorical nerve constructions and the Gepner--Haugseng universal property referred to in the supplied \(2\)-categorical result be extended to higher dimensions? The intended higher categorical structures, dimensions, hypotheses, equivalence notion, and universal property remain unspecified. In particular, the record does not determine whether the intended extension concerns each finite dimension \(n\geq 3\), an infinite-dimensional setting, or a particular higher-dimensional case. Thus the supplied evidence supports this research question but does not determine a fully quantified conjecture.
\subsection{Short English statement}
The question asks for higher-dimensional extensions of the nerve constructions and universal property mentioned in a supplied \(2\)-categorical theorem. The relevant models, equivalences, dimensions, and universal property must be recovered from the underlying references before a precise conjecture can be stated.
\subsection{Sources}
\begin{itemize}
\item Local UnsolvedMath registry, permanent ID \(30006032\), title ``Higher-Dimensional Categorical Nerves,'' source problem code \texttt{catalogue entry}, in the Oberwolfach Reports Open Problems collection. The supplied record gives proposed year \(2024\) and names no proposer.
\item Supplied report citation: \emph{Homotopical Algebra and Higher Structures (2025)}, Oberwolfach Reports, DOI \texttt{10.4171/owr/2024/39}, pp.~2251--2308. Supplied source URL: \url{https://doi.org/10.4171/owr/2024/39}. The background describes Oberwolfach Reports as documenting mathematical workshops and their open-problem discussions.
\item Supplied literature source: ``Oberwolfach Report 39 (2024),'' \url{https://doi.org/10.4171/owr/2024/39}. The exact supplied literature-check date is \texttt{2026-08-22}. This is supplied dated curation, not an independently verified assessment. It labels the problem open and reports that the report presents higher dimensions as future work and that no complete higher-dimensional result resolving the extracted question was identified. The supplied background explicitly requests independent verification of time-sensitive status and cited claims.
\item The original excerpt cites \([\mathrm{MOR22}]\) for the \(2\)-categorical theorem and attributes its universal property to Gepner--Haugseng. No full bibliographic entries, theorem location, or formulation of that universal property are supplied.
\end{itemize}
\end{document}
